\subsection{Relatively prime ideals}

Let A be a ring; two ideals a, b of A are called relatively prime if \(a + b = A\) . For this to be true, it is necessary and sufficient that \(a + b\) be contained in no prime ideal (Algebra, Chapter I, § 8, no. 7, Theorem 2), in other words, that no prime ideal contain both a and b. Two distinct maximal ideals are relatively prime.

If A is a principal ideal domain (Algebra, Chapter VII, § 1), for two elements \(a\) , \(b\) of A to be relatively prime, it is necessary and sufficient, by Bezout's identity (loc. cit., no. 2, Theorem 1), that the ideals \(Aa\) and \(Ab\) be relatively prime.

PROPOSITION 3. Let a and b be two relatively prime ideals of a ring A. Let a' and b' be two ideals of A such that every element of a (resp. b) has a power in a' (resp. b'). Then a' and b' are relatively prime.

Under the given hypothesis, every prime ideal which contains \(a'\) contains a and every prime ideal which contains \(b'\) contains b. If a prime ideal contains \(a'\) and \(b'\) , then it contains a and b, which is absurd, since a and b are relatively prime; hence \(a'\) and \(b'\) are relatively prime.

PROPOSITION 4. Let \(\mathfrak{a}, \mathfrak{b}_1, \ldots, \mathfrak{b}_n\) be ideals of a ring \(\mathbf{A}\) . If \(\mathfrak{a}\) is relatively prime to each of the \(\mathfrak{b}_i (1 \leqslant i \leqslant n)\) , it is relatively prime to \(\mathfrak{b}_1\mathfrak{b}_2 \ldots \mathfrak{b}_n\) .

Let \(\mathfrak{p}\) be a prime ideal of \(\mathbf{A}\) . If \(\mathfrak{p}\) contains \(\mathfrak{a}\) and \(\mathfrak{b}_1\mathfrak{b}_2\ldots \mathfrak{b}_n\) , it contains one of the \(\mathfrak{b}_i\) (no. 1, Proposition 1), which is absurd since \(\mathfrak{a}\) and \(\mathfrak{b}_i\) are relatively prime.

PROPOSITION 5. Let \((a_{i})_{i\in I}\) be a non-empty finite family of ideals of a ring A. The following properties are equivalent:

\begin{enumerate}
\def\labelenumi{(\alph{enumi})}
\item
  For \(i \neq j\) , \(\mathfrak{a}_i\) and \(\mathfrak{a}_j\) are relatively prime.
\item
  The canonical homomorphism \(\phi: \mathbf{A} \to \prod_{i \neq j} (\mathbf{A} / a_i)\) (Algebra, Chapter II, § 1, no. 7) is surjective.
\end{enumerate}

If these hold, the intersection \(\mathfrak{a}\) of the \(\mathfrak{a}_i\) is equal to their product and the canonical homomorphism \(\psi: \mathbf{A} / \mathfrak{a} \to \prod_{i \in I} (\mathbf{A} / \mathfrak{a}_i)\) (Algebra, Chapter II, § 1, no. 7) is bijective.

We argue by induction on n the number of elements in I, the case n = 1 being trivial. Consider first the case n = 2. Then the equivalence of (a) and (b) follows from the exactness of the sequence

\[
0 \longrightarrow \mathrm{A} / (a _ {1} \cap a _ {2}) \stackrel {\psi} {\longrightarrow} (\mathrm{A} / a _ {1}) \oplus (\mathrm{A} / a _ {2}) \longrightarrow \mathrm{A} / (a _ {1} + a _ {2}) \longrightarrow 0
\]

(Algebra, Chapter II, § 1, no. 7, formula (30)). Moreover, there exist \(e_1 \in a_1\) and \(e_2 \in a_2\) such that \(1 = e_1 + e_2\) ; then, for all \(x \in a = a_1 \cap a_2\) , \(x = xe_1 + xe_2\) ; but by definition \(xe_1 \in a_1a_2\) and \(xe_2 \in a_1a_2\) , hence \(x \in a_1a_2\) ; whence \(a \subset a_1a_2\) and the converse inclusion is obvious.

In the general case, suppose that condition (a) holds and let k be an element of I and \(b_{k} = \bigcap_{i \neq k} a_{i}\) ; the induction hypothesis implies that \(b_{k} = \prod_{i \neq k} a_{i}\) and it follows from Proposition 4 that \(a_{k}\) and \(b_{k}\) are relatively prime; then

\[
a = \bigcap_ {i \in I} a _ {i} = a _ {k} \cap b _ {k} = a _ {k} b _ {k} = \prod_ {i \in I} a _ {i}
\]

by the first part of the argument and for the same reason the canonical homomorphism \(\mathrm{A}/\mathfrak{a}\to(\mathrm{A}/\mathfrak{a}_{k})\times(\mathrm{A}/\mathfrak{b}_{k})\) is bijective; by the induction hypothesis the canonical homomorphism \(\mathrm{A}/\mathfrak{b}_{k}\to\prod_{i\neq k}(\mathrm{A}/\mathfrak{a}_{i})\) is bijective and so then is the composite homomorphism

\[
\mathrm{A} / \mathfrak {a} \rightarrow (\mathrm{A} / \mathfrak {a} _ {k}) \times (\mathrm{A} / \mathfrak {b} _ {k}) \rightarrow (\mathrm{A} / \mathfrak {a} _ {k}) \times \prod_ {i \neq k} (\mathrm{A} / \mathfrak {a} _ {i}) = \prod_ {i \in \mathrm{I}} (\mathrm{A} / \mathfrak {a} _ {i})
\]

which is precisely \(\psi\) ; that is, (b) holds. Conversely, suppose that (b) holds. We show that the \(a_i\) are necessarily relatively prime in pairs. In the contrary case, there would exist an ideal \(c \neq A\) containing \(a_i\) and \(a_j\) for \(i \neq j\) . We set \(a_h' = a_h\) for \(h\) not equal to \(i\) or \(j\) and \(a_i' = a_j' = c\) ; the canonical homomorphism \(\phi': A \to \prod_{i \in I} (A / a_i')\) can be written as the composite mapping

\[
\mathrm{A} \xrightarrow {\phi} \prod_ {i \in \mathrm{I}} (\mathrm{A} / \mathfrak {a} _ {i}) \xrightarrow {f} \prod_ {i \in \mathrm{I}} (\mathrm{A} / \mathfrak {a} _ {i} ^ {\prime})
\]

\(f\) being the product of the canonical homomorphisms \(\mathbf{A} / \mathfrak{a}_i\to \mathbf{A} / \mathfrak{a}_i'\) ; clearly \(\phi^{\prime}\) is not surjective, the projection of \(\phi^{\prime}(\mathbf{A})\) onto \((\mathbf{A} / \mathfrak{a}_i')\times (\mathbf{A} / \mathfrak{a}_i')\) being the diagonal of the product \((\mathbf{A} / \mathfrak{c})\times (\mathbf{A} / \mathfrak{c})\) , which is distinct from this product since \(\mathfrak{c}\neq \mathbf{A}\) . As \(f\) is surjective, this shows that \(\phi\) is not surjective.

PROPOSITION 6. Let \((\mathfrak{a}_i)_{i\in I}\) be a non-empty finite family of ideals of a ring A which are relatively prime in pairs; let \(\mathfrak{a}\) be the intersection of the \(\mathfrak{a}_i\) . For every A-module M, the canonical mapping \(\mathbf{M} \to \prod_{i\in I} (\mathbf{M} / \mathfrak{a}_i\mathbf{M})\) is surjective and its kernel is \(\mathfrak{aM}\) .

Clearly the canonical mapping of M to \(\prod_{i\in I}(\mathbf{M} / a_i\mathbf{M})\) is zero on aM; then, by taking quotients, it defines a homomorphism \(\lambda : \mathbf{M} / a\mathbf{M}\to \prod_{i\in I}(\mathbf{M} / a_i\mathbf{M})\) . On the other hand, by Proposition 5, the canonical homomorphism

\[
\psi \colon \mathrm{A} / \mathfrak {a} \rightarrow \prod_ {i \in \mathrm{I}} (\mathrm{A} / \mathfrak {a} _ {i})
\]

is bijective. Then so is \(1_{M} \otimes \psi: M \otimes (A/a) \to M \otimes \prod_{i \in I} (A/a_i)\) . Now \(M \otimes (A/a)\) is identified with \(M/aM\) and \(M \otimes \prod_{i \in I} (A/a_i)\) with \(\prod_{i \in I} M \otimes (A/a_i)\) , which is itself identified with \(\prod_{i \in I} (M/a_iM)\) . It is immediately verified that the above identifications transforms \(1_M \otimes \psi\) into \(\lambda\) , whence the proposition.

Example. Let K be a field, \(a_{i}\) ( \(1 \leqslant i \leqslant m\) ) distinct elements of K and, for each \(i\) , let \(g_{i}\) be a polynomial in K{[}X{]}; the principal ideal (X - \(a_{i}\) ) = m is maximal in K{[}X{]}, hence, for every system \((n_{i})_{1 \leqslant i \leqslant m}\) of \(m\) integers \(\geqslant 1\) , the ideals m \(_{i}^{n_{i}}\) are relatively prime in pairs. Then it follows from Proposition 5 that there exists a polynomial \(f \in \mathrm{K}[\mathrm{X}]\) such that \(f(\mathrm{X}) \equiv g_{i}(\mathrm{X})\) (mod. \((\mathrm{X} - a_{i})^{n_{i}}\) ) for \(1 \leqslant i \leqslant m\) , the difference of two such polynomials being divisible by \(\omega(\mathrm{X}) = \prod_{i=1}^{m} (\mathrm{X} - a_{i})^{n_{i}}\) . If all the \(n_{i}\) are taken equal to 1, we find the problem is solved explicitly by Lagrange's interpolation formula (Algebra, Chapter IV, § 2, no. 4).

